\documentclass[11pt]{letter}
\usepackage[margin=2.5cm]{geometry}
\usepackage{hyperref}

\signature{Adeline Wen \& Wei Cai\\Decentralized Computing Laboratory\\School of Engineering and Technology\\University of Washington Tacoma, Tacoma, WA, USA\\weicaics@uw.edu (corresponding author)}

\address{Adeline Wen \& Wei Cai\\Decentralized Computing Laboratory\\University of Washington Tacoma\\Tacoma, WA 98402, USA}

\begin{document}

\begin{letter}{Editorial Office\\Nature Communications}

\opening{Dear Editors,}

I am writing to submit our manuscript, ``Adverse selection in tokenized carbon markets: who profited from the first pool collapse,'' for consideration as an Article in \textit{Nature Communications}.

\textbf{Why this paper matters.} Carbon markets channel billions of dollars annually toward climate mitigation, yet fewer than 16\% of credits from investigated projects represent genuine emission reductions. Our manuscript presents the first credit-level forensic reconstruction of a carbon market collapse, analysing every deposit (1,187), every withdrawal (35,432), and every participant account (28,897) in the largest tokenized carbon credit pool (22 million tonnes). The findings directly contradict the prevailing industry and academic narrative: the pool collapsed because of CDM-era renewable energy credits with near-zero additionality (69.1\% of content), not REDD+ forest credits (4.2\%) as widely claimed.

\textbf{Key contributions.} (1) We identify the specific actors who profited: five accounts extracted 1.55 million tonnes of high-demand credits at an estimated profit of \$4--8 million, while 9.6 million tonnes of low-quality credits remain permanently stranded. (2) A natural experiment comparing identical credits across two pools with different quality screens shows that the same credits were 100\% extracted from the unscreened pool but only 28.5\% from the quality-gated pool, establishing that market design rather than credit quality drives asset fate. (3) Cross-domain evidence from two additional asset classes (staking derivatives and NFT vaults) confirms that uniform pricing of heterogeneous assets produces selective extraction across different market settings.

\textbf{Broader significance.} The Article 6.4 crediting mechanism under the Paris Agreement is currently designing pooled systems at sovereign scale. Our finding that transparency alone does not prevent quality collapse without quality-differentiated pricing speaks directly to these negotiations and to the Integrity Council for the Voluntary Carbon Market's Core Carbon Principles framework. The manuscript engages a readership spanning environmental economics, climate policy, and market design.

\textbf{Suggested reviewer expertise areas:}
\begin{itemize}
\item Environmental economics / carbon credit quality (e.g., researchers working on VCM integrity, additionality assessment)
\item Market design / mechanism design theory (e.g., researchers working on adverse selection, pooling equilibria)
\item Blockchain / tokenized asset empirics (e.g., researchers working on DeFi market microstructure)
\end{itemize}

This manuscript has not been published or submitted elsewhere. All data, analysis scripts, and scoring rubrics are publicly available at \url{https://github.com/Adeline117/carbon-neutrality} under an MIT licence, enabling full reproducibility. The authors declare no competing interests.

\closing{Sincerely,}

\end{letter}
\end{document}
